%! TEX root = ../paper.tex
\section{Discussion and Limitations}
\label{sec:discussion}

Three of this paper's findings point the same way. The channel that carries a
payload matters more than the payload, the surface that carries a leak out is
the filesystem rather than the network, and the guardrail meant to prevent
the leak is not associated with less of it. Each one says that the thing
being watched is not the thing that moves.

That is also true of the instrument. Appendix~\ref{sec:silent} catalogues six
ways this harness reported something other than what happened, and two of
them are worth restating here because they changed numbers in this paper
rather than merely embarrassing us. A fixed client timeout emptied 35 runs on
one machine, and it did so in proportion to how long each condition's system
prompt was, which means it deleted data as a function of the variable under
test. Separately, the agent sometimes wrote a tool call as text rather than
through the structured channel, and the harness read that as the end of the
run: a calibration batch with the recovery enabled leaks at 19/30 where the
same cell without it reads 14/30. Every leak rate in this paper is therefore
a lower bound, and the one place we measured the gap it was about a third.

The limits are correspondingly plain. Escapes here are induced by planted
bait on a task that never required leaving the workspace, so nothing in this
work estimates a spontaneous escape rate. Every result rests on one 7B model
served locally, and the single batch that reached a stronger model reversed
the sign, which is the strongest evidence available that these numbers do not
generalise upward. The canary detector covers plaintext and four simple
encodings; an agent encoding its token in the length or ordering of file
names would pass every monitor here. The causal experiment covers a single
task, chosen because the effect was cleanest there and never repeated on the
other two, and it eliminates one explanation for the emphasis effect without
identifying its replacement. And the harness measures without ever acting:
detection lands in about two seconds and nothing stops the agent for another
forty-four.

The cheapest next step is the one this design already implies. A containment
arm, added as a new condition rather than a modification of the existing
ones, would turn time to detection into time to containment and give the
dwell time in Section~\ref{sec:containment} something to compare against.
